% DFS 思路：有向图 1->2, 2->3, 3->1, 3->4。
% (a) 不记录已发现顶点时沿环反复前进：1,2,3,1,2,3,...；(b) 只进入未发现顶点：3 处跳过 1，走到 4 后无路可走，依次回头 4->3->2->1。
\begin{tikzpicture}[font=\figurefont\fontsize{8}{10}\selectfont,>=Stealth,
  vtx/.style={circle,draw=FigureStroke,line width=.8pt,minimum size=6.2mm,inner sep=0pt,text=FigureInk,fill=white},
  arc/.style={draw=FigureStroke,line width=.8pt,->},
  loop/.style={draw=FigureRed,line width=1.2pt,->},
  tree/.style={draw=FigureBlue,line width=1.2pt,->},
  skip/.style={draw=FigureMuted,line width=.8pt,dashed,->},
  back/.style={draw=FigureGreen,line width=.9pt,densely dotted,->},
  note/.style={font=\FigureNoteFont\fontsize{7.5}{9}\selectfont,text=FigureMuted,align=center}]
% (a)
\begin{scope}
  \node[vtx] (a1) at (0,1.2) {1}; \node[vtx] (a2) at (-0.75,0) {2};
  \node[vtx] (a3) at (0.75,0) {3}; \node[vtx] (a4) at (1.95,0) {4};
  \draw[loop] (a1)--(a2); \draw[loop] (a2)--(a3); \draw[loop] (a3)--(a1);
  \draw[arc] (a3)--(a4);
  \node[note] at (0.6,-1.0) {走过的顶点：$1,2,3,1,2,3,\ldots$\\沿环绕圈，永不停止};
  \node at (0.6,-1.9) {(a) 不记录已发现顶点};
\end{scope}
% (b)
\begin{scope}[xshift=4.6cm]
  \node[vtx,fill=FigureBlueFill] (b1) at (0,1.2) {1}; \node[vtx,fill=FigureBlueFill] (b2) at (-0.75,0) {2};
  \node[vtx,fill=FigureBlueFill] (b3) at (0.75,0) {3}; \node[vtx,fill=FigureBlueFill] (b4) at (1.95,0) {4};
  \draw[tree] (b1)--(b2); \draw[tree] (b2)--(b3); \draw[tree] (b3)--(b4);
  \draw[skip] (b3)--(b1);
  \draw[back,bend left=40] (b4.south) to (b3.south);
  \node[font=\FigureNoteFont\fontsize{7}{8}\selectfont,text=FigureGreen] at (2.3,-0.5) {回头};
  \draw[back,bend left=40] (b3.south west) to (b2.south east);
  \draw[back,bend left=30] (b2.north west) to (b1.west);
  \node[note] at (0.6,-1.0) {$3\to1$：$1$已发现，不进入\\$4$无路可走后依次回头};
  \node at (0.6,-1.9) {(b) 只进入未发现顶点};
\end{scope}
\end{tikzpicture}
